\justifying

Con todo mi cariño, quisiera agradecer: 

En primer lugar, a mi director Ramiro Avalos Ribas. Por ser mi mentor y referente, además de tutor informal y apoyo moral estos cuatro años que he tenido el placer de compartir el espacio académico. Su paciencia, consejos, y forma de resolver los problemas son los valores más importantes que me llevo de la carrera. 

Con mucho respeto, a mis directores de beca de investigación Alejandro José Uriz y Jorge Castiñeira Moreira. Por su constante disposición, por compartir mis alegrías y mis indignaciones, por su trabajo para darme las herramientas que facilitaron mi camino.    

A mi codirector Brian Maximiliano Gluzman, por compartir mi terquedad en momentos donde hubo que resolver un problema. Por enseñarme a soldar.

A Esteban Lucio González, quien me introdujo al Laboratorio y a la investigación, por confiar en mi potencial. Espero poder retribuir esa confianza dando todo de mí para ser una profesional íntegra y comprometida. 

A mis compañeros del Laboratorio de Comunicaciones, que nunca escatiman en brindar su apoyo sin importar el tipo de problema con el que me acerque, académico o extraacadémico. Gracias a ellos encontré un espacio de aprendizaje y crecimiento para dar mis primeros pasos en el ámbito de la investigación, pero sobre todo, un lugar de pertenencia. Entre ellos, quiero incluir a Leonardo David Vazquez, cuyo trabajo final sentó las bases para el desarrollo de este. Su emoción por el proyecto fue contagiosa, y estuvo presente (aún estando lejos) respondiendo dudas, llevando seguimiento y dando aliento en todo momento. 

A Gustavo Arenas, por permitirme pasar la mayor parte de mi trayectoria como estudiante siendo ayudante en la asignatura  Fisica 2/B. La docencia es una parte muy importante de quien soy, y estoy muy agradecida de ser y haber sido parte de una de las cátedras más lindas de mi recorrido como estudiante. 

A mis compañeros y docentes. Los temas pesados se hacen llevaderos cuando el ambiente de estudio es bueno. Porque tenemos la gran suerte, con casi todos los docentes en electrónica, de poder tocarles la puerta y que nos reciban para despejar cualquier duda, aunque no competa su materia.

A mis amigos, porque estoy segura de que sin ellos no es posible sobrellevar el nivel de exigencia que demanda una carrera en ingeniería. Por las largas horas de estudio, rondas de mates, charlas interminables soñando con el día en que nos recibamos, juegos de truco en el comedor entre clase y clase, "amiga te juro que vos podés", apoyo en momentos de felicidad y otros de mucho estrés y dificultades, personales y académicas. Algunas de estas personas se han convertido en mi familia, y es imposible imaginar mi vida sin ellas. 

A mi familia, gracias por soportar mis momentos de ansiedad y estrés, por su apoyo constante y fé ciega en mí.

Finalmente, expreso mi reconocimiento a la Universidad Nacional de Mar del Plata y en especial a la Facultad de Ingeniería. A lo largo de estos años se ha convertido en más que un espacio de formación, en un lugar de pertenencia del cual me enorgullezco enormemente. Sin instituciones públicas de calidad, muchos no tendríamos hoy la oportunidad de estudiar una carrera universitaria. 

\vspace{18pt}
Victoria C. Torres

5/9/2026

\thispagestyle{empty}